 
\documentclass[a4paper,12pt]{article}
\usepackage[ngerman]{babel}
\usepackage[utf8]{inputenc}
\usepackage{amsmath,amssymb}
\usepackage{tikz}
\usetikzlibrary{automata,positioning,arrows.meta}

\tikzset{
    ->,>=Stealth,
    node distance=2.5cm,
    every state/.style={thick,minimum size=1cm},
    initial text={},
}

\title{ETI -- Aufgabenblatt 2 \\ Hausaufgaben}
\author{}
\date{}

\begin{document}
\raggedright
\setlength{\parindent}{0pt}

{\LARGE\bfseries ETI -- Aufgabenblatt 2 \\ Hausaufgaben\par}
\vspace{2em}

\section*{Aufgabe 2.5}

Wir betrachten die Sprache
\[
B = \{a,b\}^* \circ \{a\} \circ \{a,b\}^2.
\]

\begin{enumerate}
    \item Geben Sie einen NFA $M_1$ mit 4 Zuständen an, der $B$ erkennt.
    \item Geben Sie einen DFA $M_2$ mit 8 Zuständen an, der $B$ erkennt.
\end{enumerate}

Begründen Sie jeweils kurz die Korrektheit.

\section*{Aufgabe 2.6}

Beweisen Sie, dass es keinen DFA mit 7 Zuständen gibt, der die Sprache $B$ aus der letzten Aufgabe erkennt.

\medskip
\textbf{Hinweis:} Beweisen Sie per Widerspruch. Betrachten Sie die Zustände, die ein solcher DFA nach Einlesen der Wörter $w \in \{a,b\}^3$ erreicht (Stichwort: Schubfachprinzip.) Was folgt dann für die Wörter $w \circ a$ bzw. $w \circ aa$?

\section*{Aufgabe 2.7}

Bestimmen Sie die Sprache $L$ über dem Alphabet $\Sigma = \{a,b,c\}$, die der folgende NFA akzeptiert. Geben Sie zudem eine kurze Begründung an.

\begin{flushleft}
\begin{tikzpicture}[node distance=2cm and 2.3cm, every edge/.append style={font=\small},
    lbl/.style={fill=white, inner sep=2pt}]
    \node[state, initial left] (q0) {$q_0$};
    \node[state, above=2.3cm of q0] (a1) {$a_1$};
    \node[state, above right=1.6cm and 1.7cm of a1] (a2) {$a_2$};
    \node[state, accepting, right=2.6cm of a2] (a3) {$a_3$};
    \node[state, right=2.6cm of q0] (b1) {$b_1$};
    \node[state, right=2.6cm of b1] (b2) {$b_2$};
    \node[state, accepting, right=2.6cm of b2] (b3) {$b_3$};
    \node[state, accepting, below right=1.6cm and 1.5cm of q0] (c1) {$c_1$};

    \path (q0) edge[bend left=25] node[lbl, left] {$c$} (a1)
          (a1) edge[bend left=25] node[lbl, right] {$\varepsilon$} (q0)
          (a1) edge node[lbl] {$a$} (a2)
          (a2) edge node[lbl] {$b$} (a3)
          (a3) edge[bend left=30] node[lbl] {$\varepsilon$} (a1);

    \path (q0) edge node[lbl] {$\varepsilon$} (b1)
          (b1) edge[loop above] node {$a,b,c$} ()
          (b1) edge node[lbl] {$b$} (b2)
          (b2) edge[loop above] node {$b$} ()
          (b2) edge node[lbl] {$a,b,c$} (b3)
          (b3) edge[bend left=25] node[lbl, below] {$a,c$} (b1);

    \path (q0) edge[bend left=20] node[lbl] {$a$} (c1)
          (c1) edge[bend left=20] node[lbl] {$\varepsilon$} (q0);
\end{tikzpicture}
\end{flushleft}

\section*{Aufgabe 2.8}

Aus einem Wort $w$ erhält man das umgekehrte Wort $w^R$, indem man die Zeichen von $w$ in umgekehrter Reihenfolge anordnet. Für jede Sprache $L$ definieren wir die umgekehrte Sprache $L^R$ als
\[
L^R = \{w^R \mid w \in L\}.
\]
Beweisen Sie: $L$ ist genau dann eine reguläre Sprache, wenn $L^R$ eine reguläre Sprache ist.

\end{document}
